% Generated draft: requires mathematical and source review.
\documentclass[11pt]{article}
\usepackage[margin=25mm]{geometry}
\usepackage{fontspec}
\setmainfont{Latin Modern Roman}
\usepackage{amsmath,amssymb,mathtools,unicode-math}
\setmathfont{Latin Modern Math}
\usepackage{enumitem,needspace,xurl,hyperref}
\setlength{\parindent}{0pt}
\setlength{\parskip}{5pt}
\setlength{\emergencystretch}{3em}
\newcommand{\N}{\mathbb{N}}
\newcommand{\Z}{\mathbb{Z}}
\newcommand{\Q}{\mathbb{Q}}
\newcommand{\R}{\mathbb{R}}
\newcommand{\C}{\mathbb{C}}
\newcommand{\F}{\mathbb{F}}
\newcommand{\A}{\mathbb{A}}
\newcommand{\PP}{\mathbb{P}}
\newcommand{\E}{\mathbb{E}}
\newcommand{\eps}{\varepsilon}
\newcommand{\dd}{\,\mathrm{d}}
\begin{document}
\section*{UnsolvedMath ID: 30003189}
\textbf{Source problem code:} catalogue entry
\subsection{Definitions and mathematical statement}
Let \(G\) be a finitely presented group, and let \(n\geq 1\) be an integer. Define its lower central series by \(\gamma_1(G)=G\) and \(\gamma_{m+1}(G)=[\gamma_m(G),G]\), where \([H,G]\) denotes the subgroup generated by commutators \(h^{-1}g^{-1}hg\). For the coefficient rings \(k=\Z,\Q\), or \(\F_p\) with \(p\) prime, let
\[
\varepsilon_k:kG\longrightarrow k,\qquad
\varepsilon_k\!\left(\sum_{g\in G}a_g g\right)=\sum_{g\in G}a_g,
\qquad \Delta(kG)=\ker\varepsilon_k.
\]
The sums defining elements of \(kG\) have finite support. Define the dimension subgroup
\[
D_n^k(G)=\gamma_{n,k}(G)=\{g\in G:g-1\in\Delta(kG)^n\}.
\]
For a subgroup \(H\leq G\), write \(H^r=\langle h^r:h\in H\rangle\) for \(r\geq1\). The corrected source formulation supplies the identities
\[
D_n^{\F_p}(G)
=\left\langle\gamma_m(G)^{p^j}:m\geq1,\ j\geq0,\ mp^j\geq n\right\rangle,
\qquad
D_n^{\Q}(G)=\sqrt{\gamma_n(G)},
\]
where
\[
\sqrt{\gamma_n(G)}=\{g\in G:g^r\in\gamma_n(G)\text{ for some integer }r\geq1\}.
\]
The generated subgroup in the first identity is the meaning of the source's product \(\prod_{mp^j\geq n}\gamma_m(G)^{p^j}\).

The supplied record asks the following questions.
\begin{enumerate}
\item Given a finite presentation of \(G\), an integer \(n\geq1\), and a prime \(p\), how can one efficiently compute \(G/D_n^{\F_p}(G)\)? Given a finite presentation of \(G\) and \(n\), how can one efficiently compute \(D_n^{\Z}(G)\)? The source does not specify an output representation for these computations or a complexity criterion for efficiency.
\item Does every finitely presented group \(G\) satisfy, for every \(n\geq1\),
\[
D_n^{\Z}(G)=\bigcap_{p\text{ prime}}D_n^{\F_p}(G)?
\]
\item Is \(2\) a universal exponent bound for the dimension quotients \(D_n^{\Z}(G)/\gamma_n(G)\)? In explicit quantifiers, the exponent-two interpretation is
\[
\text{for every finitely presented }G,\quad
\text{every }n\geq1,\quad
\text{and every }x\in D_n^{\Z}(G),\qquad x^2\in\gamma_n(G).
\]
The record writes this question as \(c_n=2\), where \(c_n\) bounds the exponent of \(D_n^{\Z}(G)/\gamma_n(G)\). It does not define \(c_n\) precisely. The displayed formulation treats \(2\) as an annihilating exponent, permitting trivial quotients; it does not require every quotient to have exponent exactly \(2\), or assert that the least universal bound equals \(2\) for every \(n\).
\end{enumerate}

There is a supplied textual discrepancy: the original prose states \(D_n^{\Q}(G)=\gamma_n(G)\), whereas the clean and corrected statements give the isolator formula above. This draft follows the explicitly designated corrected formulation and retains that discrepancy for review. The record labels the statement \texttt{corrected\_verified} and the problem partially solved; those are supplied curation labels, not independent verification by this draft.
\subsection{Short English statement}
For finitely presented groups, seek efficient methods to compute dimension subgroups over \(\Z\) and the corresponding quotients over \(\F_p\). The record also asks whether the integral dimension subgroup equals the intersection of its counterparts over all prime fields, and whether every integral dimension quotient has exponent dividing \(2\). The meaning of efficiency and the precise definition of \(c_n\) remain unspecified; the supplied rational formulas also require reconciliation.
\subsection{Sources}
\begin{itemize}
\item Supplied local UnsolvedMath registry record, permanent ID \(30003189\), title \emph{Algorithms and Exponent Bounds for Dimension Subgroups}, source problem code \texttt{catalogue entry}, in the Oberwolfach Reports Open Problems collection. No individual proposer is supplied; the proposed year is recorded as \(2016\).
\item Supplied report citation: \emph{Computational Group Theory (2017)}, Oberwolfach Reports, pp.~2123--2169, DOI \texttt{10.4171/owr/2016/37}: \url{https://doi.org/10.4171/owr/2016/37}. The background describes these reports as records of mathematical workshops and their open-problem discussions. The citation's \(2017\) label and the \(2016\) provenance indicators are retained without reconciliation.
\item Supplied literature source: Bartholdi--Mikhailov, dimension quotients, \url{https://doi.org/10.1112/topo.12301}.
\item Supplied literature source: Oberwolfach Report 37, \url{https://doi.org/10.4171/owr/2016/37}.
\item Literature assessment dated exactly \texttt{2026-08-21}: the supplied curation labels the problem partially solved and reports that dimension-subgroup exponents, counterexamples, and computational methods are known, while the extracted equality and universal exponent-\(2\) conjecture are not settled in full generality. This is supplied dated literature triage derived from the OWR curation, not an independently verified assessment; its time-sensitive status and cited claims require independent checking.
\item The supplied statement-verification field reports that the official problem-session summary was checked and that the reconstruction restores the product formula, rational isolator, quotient questions, prime intersection, and exponent-two conjecture. This draft records that claim without claiming to have visited or verified the sources.
\end{itemize}
\end{document}
